% Generated draft: requires mathematical and source review.
\documentclass[11pt]{article}
\usepackage[margin=25mm]{geometry}
\usepackage{fontspec}
\setmainfont{Latin Modern Roman}
\usepackage{amsmath,amssymb,mathtools,unicode-math}
\setmathfont{Latin Modern Math}
\usepackage{enumitem,needspace,xurl,hyperref}
\setlength{\parindent}{0pt}
\setlength{\parskip}{5pt}
\setlength{\emergencystretch}{3em}
\newcommand{\N}{\mathbb{N}}
\newcommand{\Z}{\mathbb{Z}}
\newcommand{\Q}{\mathbb{Q}}
\newcommand{\R}{\mathbb{R}}
\newcommand{\C}{\mathbb{C}}
\newcommand{\F}{\mathbb{F}}
\newcommand{\A}{\mathbb{A}}
\newcommand{\PP}{\mathbb{P}}
\newcommand{\E}{\mathbb{E}}
\newcommand{\eps}{\varepsilon}
\newcommand{\dd}{\,\mathrm{d}}
\begin{document}
\section*{UnsolvedMath ID: 20003188}
\textbf{Source problem code:} catalogue entry
\subsection{Definitions and mathematical statement}
Question 4.5 of Marc Burger starts from a symplectic vector space $V$ over a field $k$ and the stated central extension
\[
 1\longrightarrow W(k)\longrightarrow\widetilde{\operatorname{Sp}}(V)
       \longrightarrow\operatorname{Sp}(V)\longrightarrow1.
\]
In characteristic different from two, $W(k)$ is the Witt group of nonsingular symmetric forms modulo hyperbolic forms. For a closed oriented surface $S$ and a representation $\rho:\pi_1S\to\operatorname{Sp}(V)$, pull back the extension class and evaluate it on $[S]$ to obtain $\tau(\rho)\in W(k)$. The source asks to study this class for the tautological representation associated with an algebraic parameter family. Burger's Comment 4.6 says the point-parameter case is the classical Toledo invariant, with its numerical normalization understood.

The supplied partial result makes the family precise over the reals. Let $A$ be an integral real parameter variety, let $u:A\to\operatorname{Hom}(\pi_1S,\operatorname{Sp}(2n,\R))$ be regular, and set $K=\R(A)$. Evaluation at the generic point gives $\rho_K:\pi_1S\to\operatorname{Sp}(2n,K)$. Integrality and regularity matter: an arbitrary set map, a reducible algebraic set without a chosen component, or merely a map to a character quotient does not supply this tautological representation as stated.

Fix a finite bar cycle $z=\sum_j e_j[\gamma_j\mid\delta_j]$ representing the surface fundamental class, with $e_j\in\Z$, and a fixed Kashiwara/Maslov cocycle convention for the extension. The aid asserts that on a dense principal open $U\subset A$ one can simultaneously split the radicals of the finitely many cocycle forms with constant rank and obtain a model
\[
 \widetilde\tau_U\in W(\mathcal O(U))
 \quad\text{mapping to }\tau(\rho_K)\in W(K).
\]
For every $a\in U(\R)$ its evaluation is $\tau(\rho_a)\in W(\R)$. Its signature is constant on connected components of $U(\R)$; the corresponding normalized classical Toledo number is also constant along continuous real families, including across walls where a chosen cocycle presentation degenerates. This does not define unrestricted specialization of every class in $W(K)$ at every point of $A$.

In rank one, write $c(g)=g_{21}$ for $g\in SL_2(\R)$. The supplied cocycle convention represents a summand by the rank-one form with coefficient
\[
 F(g,h)=-c(g)c(gh)c(h),
\]
when it is nonzero, and by zero when degenerate. For a $C^1$ path $\rho_t$ put $F_j(t)=F(\rho_t(\gamma_j),\rho_t(\delta_j))$. Suppose every $F_j$ vanishing at $t_0$ has a simple zero there. Local constancy of the total signature then gives the simultaneous-wall identity
\[
 \sum_{j:F_j(t_0)=0}e_j\,\operatorname{sgn}(F'_j(t_0))=0.
\]
The jump of each such sign is twice its derivative sign, whereas the nonzero terms keep their signs. The formula requires the whole bar-cycle sum and does not assert cancellation term by term. Extension-class signs and factors must be fixed before identifying its signature with a particular classical Toledo normalization.

These specialize and constrain the invariant in the stated real integral-family setting; they do not classify it for all fields, singular parameter spaces or arbitrary families. In particular characteristic two and relative/punctured-surface versions require additional definitions.
\subsection{Short English statement}
Study the Witt-valued Toledo invariant of an algebraic family of symplectic surface-group representations. The supplied aid constructs a model on a dense open set for regular real integral families and matches its real fibers. In rank one, local constancy forces an explicit cancellation among all simple discriminant walls in a fixed surface-cycle representative.
\subsection{Sources}
\begin{itemize}
\item Local UnsolvedMath registry, permanent ID $20003188$, source code catalogue entry; canonical location \texttt{aim-topology-notes.json}, \texttt{notes[275]}.
\item American Institute of Mathematics, \emph{Representations of surface groups}, Question 4.5 and Comment 4.6, both Marc Burger, \url{https://aimath.org/WWN/surfacegroups/surfacegroups.pdf}. No separate literature references are supplied.
\item Attached research attempt 1 and summary describe generic specialization and the rank-one wall identity without full external citations. They are marked partial result, proved, and low novelty confidence, with an explicit machine-generated, non-peer-reviewed research-aid caveat requiring independent verification before citation.
\end{itemize}
\end{document}
